\section{Experimentation harness and hermetic sandbox}\label{sec:harnais-sandbox}

This part describes the recent work that unified the life cycle of calibration
harnesses and added a hermetic sandbox executable outside the laboratory. It
separates three objects: the anatomy of a harness, the shared core that builds its
execution object, and the sandbox that executes a deterministic arm under an audit
guard. Three commits trace it on the \code{dev/1} branch: \code{11b820c} (sandbox),
\code{02320d5} (closure of the guard's gaps) and \code{5101523} (shared core).

\subsection{Anatomy of a harness}

A harness is a set of files with a single responsibility, under
\code{benchmarks/experiments/}. The MQTT harness gives its complete form. The
\textbf{contract} is \code{mqtt\_reference.py}: it validates the briefing, binds
the tool to the execution boundary, provides the \textbf{D1 adapter} through the
deterministic arm \code{run\_rules} and recounts the run through
\code{evaluate\_trial}. The \textbf{LLM arm} is \code{mqtt\_reference\_llm.py};
\code{mqtt\_formats\_llm.py} is its H2 variant, which receives no parser. The
\textbf{live controller} is \code{mqtt\_reference\_live.py}, not to be confused
with the arm: it builds the LLM agent, the wrapper of a \code{RecordingClient}
that journals every transport call, and writes \code{decision.json},
\code{cost\_summary.json} and \code{run\_meta.json}. The \textbf{deployment} is
\code{mqtt\_reference\_proxmox.py}: the class \code{Proxmox} submits operations
via \code{ansible-playbook}, the class \code{Fixture} starts the broker from the
allocation.

The same anatomy is repeated for SSH, the API and the V0 control, with one
important difference: the H2 harness (\code{mqtt\_formats.py}) has neither a
Proxmox deployment nor an active live client in the sandbox, since its corpus is
frozen.

\subsection{The shared core}

The core exists for a precise reason, recorded in its header: the governance
boundary of \code{src/agent/core/executor.py} reads a fixed set of attributes via
\code{getattr} with a default value. An omission therefore manifests only when a
tool call misbehaves. Each harness rebuilt this object by hand; the core replaces
the copies with a single \code{build\_run}.

\code{REQUIRED\_ATTRIBUTES} enumerates thirteen attributes: the run directory, the
three budgets, the archival lock, the evidence lock, the call counter, the start
clock, the limit reason, then the scope, the context, the policy and the cost
tracking. One attribute deserves a particular comment:
\code{\_evidence\_integrity\_failed}. \code{check\_execution\_limits} reads it as
its first instruction and raises \code{EvidenceWriteError} if it is true;
\code{wrap\_tool.archive} sets it when the archive cannot be written. This lock
documents a rule: lost evidence must remain visible, and the run becomes
unverifiable. A harness that forgot it would fall back to the default value and
pursue a run whose archive is holed.

The refusal is two-level: \code{build\_run} checks the missing attribute and the
absent context key before returning the object, and \code{require\_attributes}
replays the same check on an arbitrary object. A deterministic variant declares
neither model nor provider.

Coverage remains incomplete: the four \code{*\_live.py} harnesses go through a
local function of the same name \code{new\_run}, and the core's \code{new\_run}
has no direct caller; the delegation goes through \code{build\_run}.

\subsection{The sandbox}

\code{benchmarks/experiments/sandbox.py} proves, on a development workstation as
well as in continuous integration, that a harness executes without a broker,
without a model and without a laboratory. It loads the frozen corpus (eight
fixtures) and, for each dossier: validates the dossier, opens a journal, executes
the chosen arm, writes the artefacts, then recounts via the existing evaluator.
Five artefacts are produced per trial: \code{dossier.json}, \code{oracle.json},
\code{answer.json}, \code{journal.jsonl} and \code{evaluation.json}. A
\code{sandbox-report.json} aggregates the trials and the \code{guard} block.

The arm is an explicit parameter, and an unknown name raises an error. This choice
is a correction: the original version read the arm from the corpus, which does not
declare it, rendering the conventional arm unreachable and causing the report to
invent the tested policy.

\subsection{The audit guard and its limits}

The guard relies on \code{sys.addaudithook}. The \code{Guard} class remains
installed in the process, but inspects events only when armed: the
\code{active()} context manager bounds the enforcement to the run.
\code{install\_guard} installs the hook only once and then rebinds the authorised
directory and the list of violations at each run; without this rebind, a second run
would observe the first run's directory and stifle its own violations. The blocked
events cover sockets, subprocesses, the \code{exec}/\code{spawn}/\code{posix\_spawn}
family, \code{ctypes} and rename and delete operations. Writes are handled
separately: the \code{open} event is refused if the flags contain a write mode and
if the resolved path leaves the authorised directory.

The honesty of this mechanism can be read in commit \code{02320d5}. A read-only
audit broke the first version with two measured escapes: \code{os.posix\_spawn}
started a process, and \code{ctypes.CDLL(None).system()} executed native code,
while the report announced a hermetic execution without violation. The report was
therefore corrected on three points: the key \code{hermetic} is replaced by
\code{guard\_clean}, which states only the absence of a listed event; a
\code{guard} block declares the mechanism, the blocked events and, explicitly, the
blind spots; and \code{claims\_isolation} is set to \code{false}. This last
indicator is false by construction: an audit hook observes events of the Python
interpreter, not system calls. It does not see native code reached outside
\code{ctypes}, nor a write through a descriptor opened before arming, nor a system
call issued outside CPython.

This limit has a methodological consequence: \code{guard\_clean: true} means
``none of the listed events occurred'', never ``the run is isolated''.

\begin{table}[H]
\centering
\caption{What the guard covers, and what it leaves unobserved.}
\label{tab:guard-limits}
\small
\begin{tabular}{@{}p{6.2cm}p{7.6cm}@{}}
\toprule
Covered, with refusal raised & Unobserved \\
\midrule
Sockets: connection, resolution, sending, listening & Native code reached without \code{ctypes}, \code{Popen} or \code{posix\_spawn} \\
Processes: \code{subprocess}, \code{os.system}, \code{exec*}, \code{fork}, \code{spawn*}, \code{posix\_spawn*} & Writes through a descriptor opened before arming \\
\code{ctypes}: loading, symbol resolution, call & System calls issued outside the CPython interpreter \\
Rename and delete: \code{os.rename}, \code{remove}, \code{replace}, \code{unlink}, \code{shutil} & Any behaviour after disarming: the hook remains installed but inert \\
Writes via \code{open} outside the authorised directory (resolved paths) & No guarantee of process isolation \\
\bottomrule
\end{tabular}
\end{table}

\subsection{What the sandbox demonstrates}

On the eight frozen fixtures, the declared-parser arm succeeds in six trials out of
eight: the two failures are the \code{variant/raw} trials, where the declared
parser does not cover the schema and abstains on every message, hence an erroneous
conclusion. The conventional arm, which reads the keys without a declared schema,
succeeds eight out of eight. Success is not counted on the conclusion alone: the
evaluator imposes it as a correct conclusion, without contradiction and without
any erroneous verdict.

Both arms produce the same offline lesson as H2: the failure on
\code{variant/raw} does not stem from a lack of intelligence but from a boundary
of the declared parser, and a script that tries restores the coverage. The sandbox
instruments the contrast; it does not measure it on a model. No LLM arm appears
there, no pre-flight attestation is produced, and neither \code{pair\_admissible}
nor \code{mt\_admissible} emerges from it.

\subsection{Sandbox limits}

Five limits must be read against the preceding results. First, the sandbox
involves no model: its success is a property of the harness, not an experimental
result, and the report carries \code{claims\_measurement: false}. Next,
\code{guard\_clean} is a Python observation and not a proof of isolation. Third,
the scope is partial: the negative tests that prove the hook really refuses
network, processes, \code{posix\_spawn} and \code{ctypes} go through direct
invocation of the guard, and no end-to-end test confirms that the run actually
triggers the installed hook; the detection path is proved only at unit level.
Fourth, the corpus is frozen, has a single scenario and eight fixtures, which
suffices to instrument the contrast but not to estimate a rate. Finally, the
sandbox bypasses the governance boundary it is meant to accompany: it calls the
arm directly, without building a run or passing through the tool wrapper, and the
journal is written by the harness itself; the shared core and the guard remain two
separate instruments, each verified in isolation.

\subsection{Tests}

The set counts twenty-four tests: fourteen for the sandbox and ten for the core.
They prove four properties. For the core, the run exposes the governance surface,
an incomplete run is refused at construction, \code{check\_execution\_limits}
accepts a new run, and the exceeding of the call budget is raised by the tool
wrapper and not by the limit check. For the sandbox, the eight fixtures produce the
five artefacts and a consistent report, and the arm is declared and never read
from the corpus. For the H2 lesson, the declared parser fails on at least one trial
when the conventional arm succeeds everywhere. For the guard, negative tests
refuse writes outside the directory, the network, subprocesses,
\code{posix\_spawn} and \code{ctypes}, while a test verifies that the hook remains
inert after disarming and that a second guard observes its own directory.

What these tests do not prove: none executes an LLM arm, none measures a
confidence interval, and none verifies that the installed hook is reached during
the run rather than only when it is called by hand.
